\documentclass[UTF8, a4paper, 12pt]{ctexart}

% 数学宏包
\usepackage{amsmath}
\usepackage{amssymb}
\usepackage{amsthm}
\usepackage{mathtools}

% 页面设置
\usepackage[a4paper, margin=2.5cm]{geometry}

% 超链接
\usepackage[colorlinks=true, linkcolor=blue, urlcolor=blue]{hyperref}

% 范畴论图示
\usepackage{tikz-cd}

% 列表设置
\usepackage{enumitem}

% 取消首行缩进
\ctexset{autoindent=false}   % 先关闭 ctex 的自动调整功能
\setlength{\parindent}{0pt}  % 再把段首缩进设为 0

% 常用数学符号
\newcommand{\Set}{\mathbf{Set}}
\newcommand{\Grp}{\mathbf{Grp}}
\newcommand{\Ab}{\mathbf{Ab}}
\newcommand{\Cat}{\mathbf{Cat}}
\newcommand{\id}{\mathrm{id}}
\newcommand{\op}{\mathrm{op}}
\newcommand{\normalcl}[1]{\left\langle #1 \right\rangle^{\mathrm{normal}}}

\title{范畴论中的 Equalizer 与 Coequalizer}
\author{}
\date{}

\begin{document}

\maketitle

\section*{问题}

范畴论中的 equalizer、coequalizer 的定义是什么？
它们的直观是什么？请给出例子解释。

\section{Equalizer：等化子}

设在一个范畴中有两个平行态射
\[
f,g\colon X\rightrightarrows Y,
\]
也就是两个态射
\[
f\colon X\to Y,\qquad g\colon X\to Y.
\]

\subsection{定义}

如果有态射
\[
e\colon E\to X
\]
满足
\[
f\circ e=g\circ e,
\]
并且对任意态射
\[
h\colon Z\to X,
\]
只要
\[
f\circ h=g\circ h,
\]
就存在唯一的态射
\[
u\colon Z\to E
\]
使得
\[
h=e\circ u,
\]
那么称 \(e\colon E\to X\) 是 \(f,g\) 的
\emph{equalizer}，中文称为\textbf{等化子}。

其交换图如下：
\[
\begin{tikzcd}
Z \arrow[rr, "u", dashed] \arrow[dr, "h"'] &&
E \arrow[dl, "e"] \\
& X &
\end{tikzcd}
\]
其中还要求
\[
f\circ e=g\circ e.
\]

完整地说，等化子满足以下两个条件：

\begin{enumerate}[label=\arabic*.]
    \item \(f\circ e=g\circ e\)；
    \item 对任意满足 \(f\circ h=g\circ h\) 的 \(h\colon Z\to X\)，
    存在唯一的 \(u\colon Z\to E\)，使得 \(h=e\circ u\)。
\end{enumerate}

\subsection{直观理解}

Equalizer 是 \(X\) 中那些被 \(f\) 和 \(g\) ``区分不出来''的部分。

在集合范畴 \(\Set\) 中，等化子就是集合
\[
E=\{x\in X\mid f(x)=g(x)\},
\]
以及它到 \(X\) 的包含映射
\[
e\colon E\hookrightarrow X.
\]

因此，equalizer 可以直观地理解为：

\begin{quote}
取出 \(X\) 中使 \(f\) 和 \(g\) 相等的最大子对象。
\end{quote}

这里的“最大”不是简单地指元素数量最多，而是指它满足一个泛性质：
任何其他使 \(f\) 和 \(g\) 相等的对象，都必须唯一地经过 \(E\)。

\subsection{集合中的例子}

设
\[
X=\mathbb{R},\qquad Y=\mathbb{R},
\]
定义
\[
f(x)=x^2,\qquad g(x)=|x|.
\]

我们寻找满足
\[
x^2=|x|
\]
的实数 \(x\)。

解为
\[
x=-1,\qquad x=0,\qquad x=1.
\]

所以 equalizer 是集合
\[
E=\{-1,0,1\},
\]
以及包含映射
\[
e\colon E\hookrightarrow \mathbb{R}.
\]

因为对于每个 \(x\in E\)，都有
\[
f(x)=g(x).
\]

如果另一个集合 \(Z\) 有映射
\[
h\colon Z\to\mathbb{R},
\]
并且满足
\[
f\circ h=g\circ h,
\]
那么对于每个 \(z\in Z\)，都有
\[
h(z)\in\{-1,0,1\}.
\]

因此，\(h\) 唯一地分解为
\[
Z\xrightarrow{u}E\xrightarrow{e}\mathbb{R}.
\]

\section{Coequalizer：余等化子}

仍然设有两个平行态射
\[
f,g\colon X\rightrightarrows Y.
\]

\subsection{定义}

如果有态射
\[
q\colon Y\to Q
\]
满足
\[
q\circ f=q\circ g,
\]
并且对任意态射
\[
h\colon Y\to Z,
\]
只要
\[
h\circ f=h\circ g,
\]
就存在唯一的态射
\[
u\colon Q\to Z
\]
使得
\[
h=u\circ q,
\]
那么称 \(q\colon Y\to Q\) 是 \(f,g\) 的
\emph{coequalizer}，中文称为\textbf{余等化子}。

其交换图如下：
\[
\begin{tikzcd}
X
    \arrow[r, bend left=45, "f"]
    \arrow[r, bend right=45, swap, "g"]
&
Y
    \arrow[r, "q"]
    \arrow[dr, "h"']
&
Q
    \arrow[d, "u", dashed]
\\
&& Z
\end{tikzcd}
\]

其中要求
\[
q\circ f=q\circ g,
\]
并且任意满足
\[
h\circ f=h\circ g
\]
的态射 \(h\colon Y\to Z\)，都唯一地分解为
\[
h=u\circ q.
\]

\subsection{直观理解}

Coequalizer 不再从 \(X\) 中选出一部分，而是把 \(Y\) 中的某些元素识别起来，使得 \(f\) 和 \(g\) 的结果变得相同。

在集合范畴 \(\Set\) 中，coequalizer 大致是商集合
\[
Q=Y/{\sim},
\]
其中 \(\sim\) 是由关系
\[
f(x)\sim g(x)\qquad (x\in X)
\]
生成的最小等价关系。

因此，coequalizer 可以直观地理解为：

\begin{quote}
把 \(f(x)\) 和 \(g(x)\) 强制识别为同一个元素，从而得到最一般的商对象。
\end{quote}

\subsection{集合中的例子}

设
\[
X=\{a,b\},\qquad Y=\{1,2,3,4\}.
\]

定义两个映射
\[
f(a)=1,\qquad f(b)=2,
\]
以及
\[
g(a)=2,\qquad g(b)=3.
\]

我们要求一个映射
\[
q\colon Y\to Q
\]
满足
\[
q\circ f=q\circ g.
\]

由 \(a\) 得到
\[
q(1)=q(2).
\]

由 \(b\) 得到
\[
q(2)=q(3).
\]

因此必须有
\[
q(1)=q(2)=q(3).
\]

元素 \(4\) 不受任何关系约束，可以单独保留。

于是 coequalizer 可以取为
\[
Q=\{[1],[4]\},
\]
其中
\[
[1]=\{1,2,3\},\qquad [4]=\{4\}.
\]

商映射为
\[
q\colon Y\to Q,
\]
满足
\[
q(1)=q(2)=q(3)=[1],
\qquad
q(4)=[4].
\]

这就是把由 \(f(x)\) 和 \(g(x)\) 产生的元素强行合并。

\section{群中的 Equalizer}

在群范畴 \(\Grp\) 中，设有两个群同态
\[
f,g\colon G\to H.
\]

那么它们的 equalizer 是子群
\[
E=\{x\in G\mid f(x)=g(x)\}.
\]

特别地，如果 \(g\) 是平凡同态，或者考虑
\[
f\colon G\to H,\qquad 0\colon G\to H,
\]
那么 equalizer 就是 \(f\) 的核：
\[
\ker f
=
\{x\in G\mid f(x)=e_H\}.
\]

因此：

\begin{quote}
群同态的 kernel 是某个 equalizer 的特殊情形。
\end{quote}

\section{群中的 Coequalizer}

在群范畴 \(\Grp\) 中，设有两个群同态
\[
f,g\colon G\to H.
\]

那么 coequalizer 是一个商群
\[
q\colon H\to Q
\]
满足
\[
q\circ f=q\circ g.
\]

要求
\[
q(f(x))=q(g(x)),
\]
等价于要求
\[
q\bigl(f(x)g(x)^{-1}\bigr)=e_Q.
\]

因此，需要把所有元素
\[
f(x)g(x)^{-1}
\]
都变成单位元。为了保持群结构，需要取这些元素生成的正规闭包：
\[
N=
\normalcl{
f(x)g(x)^{-1}\mid x\in G
}.
\]

于是 coequalizer 是自然商映射
\[
q\colon H\to H/N.
\]

特别地，如果 \(g\) 是平凡同态，则
\[
N=\normalcl{f(G)}.
\]

如果 \(f(G)\) 本身是正规子群，则 coequalizer 就是
\[
H/f(G).
\]

这说明：

\begin{quote}
在群范畴中，余等化子通常表现为除以某个正规子群。
\end{quote}

\section{Equalizer 与 Coequalizer 的对偶性}

Equalizer 和 coequalizer 是完全对偶的概念。

Equalizer 的形式是
\[
E\longrightarrow X\rightrightarrows Y,
\]
也就是寻找一个对象 \(E\to X\)，使得两条复合路径相等。

Coequalizer 的形式是
\[
X\rightrightarrows Y\longrightarrow Q,
\]
也就是从 \(Y\) 出发，通过商映射使两条平行态射相等。

如果把范畴中的所有箭头反向：

\begin{itemize}
    \item equalizer 会变成 coequalizer；
    \item coequalizer 会变成 equalizer。
\end{itemize}

形式化地说：
\[
\operatorname{Coeq}_{\mathcal{C}}(f,g)
=
\operatorname{Eq}_{\mathcal{C}^{\op}}(f^{\op},g^{\op}).
\]

\section{与 Kernel、Cokernel 的关系}

在有零对象的范畴中：

\begin{itemize}
    \item kernel 是某种 equalizer；
    \item cokernel 是某种 coequalizer。
\end{itemize}

对于群或阿贝尔群，有
\[
\ker(f)=\operatorname{Eq}(f,0).
\]

在阿贝尔群范畴 \(\Ab\) 中，
\[
\operatorname{coker}(f)=\operatorname{Coeq}(f,0).
\]

具体地，对于阿贝尔群同态
\[
f\colon G\to H,
\]
有
\[
\operatorname{coker}(f)
=
H/\operatorname{im}(f).
\]

在一般群中，则需要对像取正规闭包。

\section{总结}

对于平行态射
\[
f,g\colon X\rightrightarrows Y,
\]
可以这样记忆：

\subsection*{Equalizer}

\[
E\to X
\]

是 \(X\) 中使 \(f\) 和 \(g\) 相等的部分。

在集合中：
\[
E=\{x\in X\mid f(x)=g(x)\}.
\]

\subsection*{Coequalizer}

\[
Y\to Q
\]

是把 \(f(x)\) 和 \(g(x)\) 识别为相同之后得到的商对象。

在集合中：
\[
Q=Y/\bigl(f(x)\sim g(x)\bigr).
\]

一句话概括：

\begin{quote}
Equalizer 是取满足相等条件的子对象；
Coequalizer 是施加相等关系得到商对象。
\end{quote}

\end{document}